\begin{tomtat}
	\begin{minipage}[t]{0.3\textwidth}
		\vspace{0pt}
		\begin{tomtat}
			\begin{center}
				\textbf{ĐỀ 53}
			\end{center}
			
		\end{tomtat}
	\end{minipage}
	\hfill
	\begin{minipage}[t]{0.5\textwidth}
		\vspace{0pt}
		\begin{tomtat}
			\begin{center}
				\includegraphics[scale=0.08]{images/53}
			\end{center}
		\end{tomtat}
	\end{minipage}
\end{tomtat}
\textbf{I. PHẦN TRẮC NGHIỆM}
\begin{ex}
	Điều kiện xác định của phương trình $\dfrac{3}{x-1} - \dfrac{2x+3}{x+2} = \dfrac{3-x}{2}$ là:
	\choice{$x \ne -1, x \ne -2$}{$x \ne -1, x \ne 2$}{$x \ne 1, x \ne -2$}{$x \ne 1, x \ne 2$}
	\loigiai{
		\textbf{Đáp án đúng là C.} \\
		Điều kiện xác định của phương trình là các mẫu số phải khác 0:
		\begin{itemize}
			\item $x - 1 \ne 0 \Leftrightarrow x \ne 1$
			\item $x + 2 \ne 0 \Leftrightarrow x \ne -2$
		\end{itemize}
		Vậy, điều kiện xác định là $x \ne 1$ và $x \ne -2$.
	}
\end{ex}

\begin{ex}
	Tính $\sqrt{16x^2y^4}$:
	\choice{$4xy^2$}{$4x^2y^4$}{$4|x|y^2$}{$-4xy^2$}
	\loigiai{
		\textbf{Đáp án đúng là C.} \\
		Ta có:
		\begin{align*}
			\sqrt{16x^2y^4} &= \sqrt{(4xy^2)^2} \\
			&= |4xy^2| \\
			&= 4|x|(y^2) \quad (\text{vì } y^2 \ge 0) \\
			&= 4|x|y^2
		\end{align*}
	}
\end{ex}

\begin{ex}
	Trong cùng một mặt phẳng tọa độ Oxy cho đồ thị của bốn hàm số sau: $y=4x^2$; $y=-3x^2$; $y=-10x^2$; $y=-4x^2$. Trong số các đồ thị đã cho, có bao nhiêu đồ thị nằm phía dưới trục hoành?
	\choice{$1$}{$3$}{$2$}{$4$}
	\loigiai{
		\textbf{Đáp án đúng là B.} \\
		Đồ thị hàm số $y=ax^2$ nằm phía dưới trục hoành khi $a < 0$.
		Các hàm số có $a < 0$ là:
		\begin{itemize}
			\item $y = -3x^2$ (có $a = -3 < 0$)
			\item $y = -10x^2$ (có $a = -10 < 0$)
			\item $y = -4x^2$ (có $a = -4 < 0$)
		\end{itemize}
		Hàm số $y=4x^2$ có $a=4>0$ nên đồ thị nằm phía trên trục hoành.
		Vậy có 3 đồ thị nằm phía dưới trục hoành.
	}
\end{ex}

\begin{ex}
	Bất phương trình $5x - 4 - 3(2x-9) \le 5x - 8$ có tập nghiệm là:
	\choice{$x \ge \dfrac{31}{6}$}{$x \ge \dfrac{-31}{6}$}{$x < \dfrac{31}{6}$}{$x > \dfrac{31}{6}$}
	\loigiai{
		\textbf{Đáp án đúng là A.} \\
		Ta có:
		\begin{align*}
			5x - 4 - 3(2x-9) &\le 5x - 8 \\
			\Leftrightarrow 5x - 4 - 6x + 27 &\le 5x - 8 \\
			\Leftrightarrow -x + 23 &\le 5x - 8 \\
			\Leftrightarrow 23 + 8 &\le 5x + x \\
			\Leftrightarrow 31 &\le 6x \\
			\Leftrightarrow x &\ge \dfrac{31}{6}
		\end{align*}
		Vậy tập nghiệm của bất phương trình là $x \ge \dfrac{31}{6}$.
	}
	\end{ex}
\begin{ex}
	Cho tam giác ABC vuông tại A có $BC = 8 \text{ cm}$, $AC = 6 \text{ cm}$. Tỉ số lượng giác $\tan C$ là bao nhiêu? (kết quả được làm tròn đến hàng phần trăm)
	\choice{$0,87$}{$0,86$}{$0,88$}{$0,89$}
	\loigiai{
		\textbf{Đáp án đúng là C} \\
		Xét tam giác ABC vuông tại A, theo định lý Py-ta-go ta có:
		\begin{align*}
			AB^2 + AC^2 &= BC^2 \\
			AB^2 + 6^2 &= 8^2 \\
			AB^2 + 36 &= 64 \\
			AB^2 &= 64 - 36 \\
			AB^2 &= 28 \\
			AB &= \sqrt{28} = 2\sqrt{7} \text{ (cm)}
		\end{align*}
		Tỉ số lượng giác $\tan C$:
		\begin{align*}
			\tan C &= \frac{\text{cạnh đối}}{\text{cạnh kề}} = \frac{AB}{AC} \\
			&= \frac{2\sqrt{7}}{6} \\
			&= \frac{\sqrt{7}}{3} \\
			&\approx \frac{2.6457513}{3} \\
			&\approx 0.881917...
		\end{align*}
		Làm tròn đến hàng phần trăm, $\tan C \approx 0,88$.
	}
\end{ex}

\begin{ex}
	Cung có số đo $110^\circ$ của đường tròn bán kính $8 \text{ cm}$ dài bao nhiêu? (làm tròn kết quả đến hàng phần chục)
	\choice{$15,3 \text{ cm}$}{$15,4 \text{ cm}$}{$15,5 \text{ cm}$}{$15 \text{ cm}$}
	\loigiai{
		\textbf{Đáp án đúng là B.} \\
		Độ dài cung tròn được tính bằng công thức: $l = \dfrac{\pi R n}{180}$, trong đó $R$ là bán kính và $n$ là số đo cung (bằng độ).
		\begin{align*}
			l &= \dfrac{\pi \times 8 \text{ cm} \times 110}{180} \\
			&= \dfrac{880\pi}{180} \text{ cm} \\
			&= \dfrac{44\pi}{9} \text{ cm} \\
			&\approx \dfrac{44 \times 3.14159265}{9} \text{ cm} \\
			&\approx \dfrac{138.2300766}{9} \text{ cm} \\
			&\approx 15.3588974... \text{ cm}
		\end{align*}
		Làm tròn đến hàng phần chục, $l \approx 15,4 \text{ cm}$.
	}
\end{ex}


\begin{ex}
	Đội tuyển kéo co của lớp 9A gồm có 6 bạn nam và 4 bạn nữ. Chọn ngẫu nhiên một bạn đứng đầu hàng. Xác suất “Bạn nam được chọn đứng đầu hàng” là:
	\choice{$\dfrac{3}{5}$}{$\dfrac{2}{5}$}{$\dfrac{2}{3}$}{$\dfrac{1}{3}$}
	\loigiai{
		\textbf{Đáp án đúng là A.} \\
		Tổng số học sinh trong đội tuyển kéo co là $6 (\text{nam}) + 4 (\text{nữ}) = 10$ (bạn).
		Số bạn nam là 6 bạn.
		Gọi A là biến cố “Bạn nam được chọn đứng đầu hàng”.
		Số kết quả thuận lợi cho biến cố A (chọn 1 bạn nam từ 6 bạn nam) là $n(A) = 6$.
		Số phần tử của không gian mẫu (chọn 1 bạn bất kỳ từ 10 bạn) là $n(\Omega) = 10$.
		Xác suất của biến cố A là:
		$$ P(A) = \frac{n(A)}{n(\Omega)} = \frac{6}{10} = \frac{3}{5} $$
	}
\end{ex}

\textbf{II. PHẦN TỰ LUẬN}
\begin{ex}
	(1 điểm) Giải hệ phương trình: $\begin{cases} x+y=3 \\ 2x-y=3 \end{cases}$
	\loigiai{
		\textbf{Đáp án:} \\
		Cộng vế theo vế hai phương trình của hệ, ta được:
		\begin{align*}
			(x+y) + (2x-y) &= 3+3 \\
			3x &= 6 \\
			x &= 2
		\end{align*}
		Thay $x=2$ vào phương trình thứ nhất $x+y=3$, ta được:
		\begin{align*}
			2+y &= 3 \\
			y &= 3-2 \\
			y &= 1
		\end{align*}
		Vậy nghiệm của hệ phương trình là $(x; y) = (2; 1)$.
	}
\end{ex}

\begin{ex}
	(1 điểm) Rút gọn biểu thức: $B = \dfrac{1}{\sqrt{x}+2} + \dfrac{\sqrt{x}}{2-\sqrt{x}} + \dfrac{2x-\sqrt{x}+2}{x-4}$ với $x > 0, x \ne 4$.
	\loigiai{
		\textbf{Đáp án:} \\
		Với $x > 0, x \ne 4$, ta có:
		\begin{align*}
			B &= \dfrac{1}{\sqrt{x}+2} + \dfrac{\sqrt{x}}{2-\sqrt{x}} + \dfrac{2x-\sqrt{x}+2}{(\sqrt{x}-2)(\sqrt{x}+2)} \\
			&= \dfrac{1}{\sqrt{x}+2} - \dfrac{\sqrt{x}}{\sqrt{x}-2} + \dfrac{2x-\sqrt{x}+2}{(\sqrt{x}-2)(\sqrt{x}+2)} \\
			&= \dfrac{(\sqrt{x}-2) - \sqrt{x}(\sqrt{x}+2) + (2x-\sqrt{x}+2)}{(\sqrt{x}-2)(\sqrt{x}+2)} \\
			&= \dfrac{\sqrt{x}-2 - x-2\sqrt{x} + 2x-\sqrt{x}+2}{(\sqrt{x}-2)(\sqrt{x}+2)} \\
			&= \dfrac{(2x-x) + (\sqrt{x}-2\sqrt{x}-\sqrt{x}) + (-2+2)}{(\sqrt{x}-2)(\sqrt{x}+2)} \\
			&= \dfrac{x - 2\sqrt{x}}{(\sqrt{x}-2)(\sqrt{x}+2)} \\
			&= \dfrac{\sqrt{x}(\sqrt{x}-2)}{(\sqrt{x}-2)(\sqrt{x}+2)} \\
			&= \dfrac{\sqrt{x}}{\sqrt{x}+2}
		\end{align*}
		Vậy $B = \dfrac{\sqrt{x}}{\sqrt{x}+2}$.
	}
\end{ex}

\begin{ex}
	(1 điểm) Cho phương trình $x^2-4x+1=0$. Không giải phương trình hãy tính giá trị biểu thức sau: $T = x_1^2 + x_2^2 - 5(x_1+x_2)$.
	\loigiai{
		\textbf{Đáp án:} \\
		Phương trình $x^2-4x+1=0$ có $a=1, b=-4, c=1$.
		Vì $\Delta' = (-2)^2 - 1 \cdot 1 = 4-1 = 3 > 0$, phương trình có hai nghiệm phân biệt $x_1, x_2$.
		Theo hệ thức Vi-ét, ta có:
		\begin{itemize}
			\item $x_1 + x_2 = -\dfrac{b}{a} = -\dfrac{-4}{1} = 4$
			\item $x_1 x_2 = \dfrac{c}{a} = \dfrac{1}{1} = 1$
		\end{itemize}
		Ta có:
		\begin{align*}
			T &= x_1^2 + x_2^2 - 5(x_1+x_2) \\
			&= (x_1+x_2)^2 - 2x_1x_2 - 5(x_1+x_2) \\
			&= (4)^2 - 2(1) - 5(4) \\
			&= 16 - 2 - 20 \\
			&= 14 - 20 \\
			&= -6
		\end{align*}
		Vậy $T = -6$.
	}
\end{ex}
\begin{ex}
	Nhóm học sinh tình nguyện khối 9 của một trường trung học cơ sở có các bạn sau:
	\begin{itemize}
		\item 2 bạn nam là: Trung (lớp 9A); Việt (lớp 9C).
		\item 3 bạn nữ là: An (lớp 9A); Châu (lớp 9B); Hương (lớp 9C).
	\end{itemize}
	Chọn ngẫu nhiên 1 bạn nam và 1 bạn nữ trong nhóm đó để tham gia hoạt động tình nguyện của trường. Xác định không gian mẫu của phép thử và tính xác suất của biến cố sau: “Hai bạn được chọn ra đều là học sinh lớp 9A”.
	\loigiai{
		\textbf{Lời giải:} \\
		Gọi các bạn nam là $N_1$=Trung(9A), $N_2$=Việt(9C).
		Gọi các bạn nữ là $NU_1$=An(9A), $NU_2$=Châu(9B), $NU_3$=Hương(9C).
		
		Phép thử là chọn ngẫu nhiên 1 bạn nam và 1 bạn nữ.
		Không gian mẫu $\Omega$ (các cặp (nam, nữ) có thể được chọn) được liệt kê như sau:
		\begin{align*}
			\Omega = \{ &(N_1, NU_1); (N_1, NU_2); (N_1, NU_3); \\
			&(N_2, NU_1); (N_2, NU_2); (N_2, NU_3) \}
		\end{align*}
		Thay tên cụ thể:
		\begin{align*}
			\Omega = \{ &(\text{Trung (9A), An (9A)}); (\text{Trung (9A), Châu (9B)}); (\text{Trung (9A), Hương (9C)}); \\
			&(\text{Việt (9C), An (9A)}); (\text{Việt (9C), Châu (9B)}); (\text{Việt (9C), Hương (9C)}) \}
		\end{align*}
		Số phần tử của không gian mẫu là $n(\Omega) = 6$.
		
		Gọi A là biến cố: “Hai bạn được chọn ra đều là học sinh lớp 9A”.
		Để biến cố A xảy ra, bạn nam được chọn phải là Trung (9A) và bạn nữ được chọn phải là An (9A).
		Các phần tử thuận lợi cho biến cố A được liệt kê là:
		$$ A = \{ (\text{Trung (9A), An (9A)}) \} $$
		Số phần tử của biến cố A là $n(A) = 1$.
		
		Xác suất của biến cố A là:
		$$ P(A) = \frac{n(A)}{n(\Omega)} = \frac{1}{6} $$
		Vậy, xác suất để hai bạn được chọn ra đều là học sinh lớp 9A là $\dfrac{1}{6}$.
		
	}
\end{ex}
\begin{ex} Bác An đi xe máy từ nhà đến nơi làm việc cách nhau 60km với vận tốc dự định trước. Sau khi đi được $\dfrac{1}{3}$ quãng đường, do điều kiện thời tiết không thuận lợi nên trên quãng đường còn lại bác An phải đi với vận tốc ít hơn so với vận tốc dự định ban đầu 10 km/h. Tính vận tốc dự định của bác An khi đi từ nhà đến nơi làm việc, biết bác An đến nơi làm việc muộn hơn dự định 20 phút.
	\loigiai{
		Gọi vận tốc dự định của bác An là $x$ (km/h, $x > 10$ vì vận tốc lúc sau là $x-10$).
		Thời gian dự định đi hết quãng đường là $\dfrac{60}{x}$ (giờ).
		
		Quãng đường đầu bác An đi được là: $\dfrac{1}{3} \times 60 = 20$ km.
		Thời gian đi hết quãng đường đầu là: $\dfrac{20}{x}$ (giờ).
		
		Quãng đường còn lại là: $60 - 20 = 40$ km.
		Vận tốc trên quãng đường còn lại là: $x - 10$ (km/h).
		Thời gian đi hết quãng đường còn lại là: $\dfrac{40}{x-10}$ (giờ).
		
		Tổng thời gian thực tế bác An đi là: $\dfrac{20}{x} + \dfrac{40}{x-10}$ (giờ).
		
		Đổi 20 phút = $\dfrac{20}{60} = \dfrac{1}{3}$ giờ.
		Theo đề bài, bác An đến muộn hơn dự định 20 phút, nên ta có phương trình:
		\begin{align*}
			\left(\dfrac{20}{x} + \dfrac{40}{x-10}\right) - \dfrac{60}{x} &= \dfrac{1}{3} \\
			\dfrac{40}{x-10} - \dfrac{40}{x} &= \dfrac{1}{3} \\
			\dfrac{40x - 40(x-10)}{x(x-10)} &= \dfrac{1}{3} \\
			\dfrac{40x - 40x + 400}{x^2 - 10x} &= \dfrac{1}{3} \\
			\dfrac{400}{x^2 - 10x} &= \dfrac{1}{3} \\
			x^2 - 10x &= 400 \times 3 \\
			x^2 - 10x &= 1200 \\
			x^2 - 10x - 1200 &= 0
		\end{align*}
		
		Giải phương trình bậc hai:
		$\Delta = (-10)^2 - 4 \times 1 \times (-1200) = 100 + 4800 = 4900$
		$\sqrt{\Delta} = \sqrt{4900} = 70$
		
		Phương trình có hai nghiệm phân biệt:
		$x_1 = \dfrac{-(-10) + 70}{2 \times 1} = \dfrac{10 + 70}{2} = \dfrac{80}{2} = 40$ (Thỏa mãn điều kiện $x > 10$)
		$x_2 = \dfrac{-(-10) - 70}{2 \times 1} = \dfrac{10 - 70}{2} = \dfrac{-60}{2} = -30$ (Không thỏa mãn điều kiện $x > 10$)
		
		Vậy, vận tốc dự định của bác An là 40 km/h.
	}
\end{ex}
\begin{ex} Cho tam giác ABC nhọn nội tiếp đường tròn (O), các đường cao AD, BE, CF cắt nhau tại H. Kẻ đường kính AQ của đường tròn (O).
	\begin{enumerate}
		\item[a)] Chứng minh 4 điểm A, F, H, E cùng thuộc một đường tròn.
		\item[b)] Chứng minh rằng $\widehat{BAD} = \widehat{CAQ}$.
		\item[c)] Vẽ OM vuông góc với BC (M thuộc BC). Tính tỉ số $\dfrac{OM}{AH}$.
	\end{enumerate}
\end{ex}

\begin{ex}
Bác Long dùng 80 m dây thép gai để rào một mảnh vườn có dạng hình chữ nhật. Bác đã tận dụng một bờ tường có sẵn để làm một cạnh của hàng rào của mảnh vườn. Tìm các kích thước của mảnh vườn có diện tích lớn nhất mà bác Long rào được bằng 80 m lưới thép gai.

\end{ex}

\setcounterref{ex}{0}
